% === VI. EXPERIMENTS AND RESULTS =============================================
\section{Experiments and Results}
\label{sec:results}

\subsection{Setup}
\label{subsec:setup}

All experiments run on the architecture of Section~\ref{sec:platform}: the
Formula AI vehicle in CARLA driven through ROS~2 Humble, the server in
synchronous mode at a \SI{0.033}{\second} fixed physics step. The identifier
collects \SI{30}{\second} of data on a \SI{50}{\hertz} timer that the server's
step rate holds to \SI{30}{\hertz} (Section~\ref{subsec:data}) and
re-identifies from the newest buffer, matching the baseline's data
budget~\cite{ontrack2025}. Ground truth is the simulator's own per-wheel
telemetry, which the identifier never reads; every figure is exported with the
CSV holding its samples, and the cross-run comparison reads those CSVs rather
than re-deriving any metric.

Two clocks appear below, and no metric mixes them. The controller node stamps
its own quantities (run duration, handover time, lap time and every tracking
error) on the ROS wall clock. The identifier and the force benchmark stamp
theirs on the simulation clock, which the server advances faster than wall time,
at a ratio between \num{1.28} and \num{1.66} across the four runs. Table~\ref{tab:comparison}
labels each block with the clock it is measured on.

\subsubsection{Two friction schedules}
\label{subsubsec:surface}
Four runs are reported, two configurations crossed with two surfaces.

On the \emph{static} surface the per-wheel coefficient is held at its configured
\num{1.0}, and the road-surface factor of \num{0.70} PhysX applies on top of it
makes the effective coefficient \num{0.700}, constant to the last digit the
simulator reports over all force-validation samples of both runs. It is
deliberately not the $\mu=\num{1.05}$ surface the tire prior was measured on
(Section~\ref{subsubsec:priorprovenance}), so an arm that cannot move off its
prior cannot score well. What it does not test is tracking: the identifier is
asked to find one number and then hold it.

The \emph{decaying} surface supplies the change. A schedule node drives the
per-wheel coefficient over the ROS wall clock through
\begin{equation}
\mu_{\mathrm{cmd}}(t)=\max\!\big(\mu_{\mathrm{floor}},\;\mu_{0}(1-\eta t)\big),
\qquad \eta=10^{-3}\,\si{\per\second},
\label{eq:decay}
\end{equation}
with $\mu_{0}=\num{1.0}$ and $\mu_{\mathrm{floor}}=\num{0.5}$, reached at
$t=\SI{500}{\second}$; the command is applied through the bridge's
\texttt{set\_tire\_friction} service while the car is driving, so the surface
changes under a moving vehicle rather than between runs. The effective axle
coefficient the plant reports therefore falls monotonically from \num{0.700} to
the \num{0.350} floor, half the grip gone over the run, and the force-validation
window opens at \num{0.637}. The linear-decay protocol is that of
LLA-MPC~\cite{llampc2025}, which decays $\mu$ at \SI{2}{\percent} of nominal per
second on a 1:43 vehicle. We keep the law and the total excursion and reduce the
rate by a factor of \num{20}, because the scaling of the platform sets what an
on-track identifier can be scored against at all: at \SI{2}{\percent} per second
the surface is exhausted in \SI{50}{\second}, shorter than one
\SI{30}{\second} buffer plus the \SI{20.9}{\second} its fit takes
(Section~\ref{subsec:runtime}), and no identifier of this class would publish a
second coefficient set before the experiment ended. At \eqref{eq:decay}'s rate
the surface instead loses about \num{0.06} of $\mu$ per identification cycle,
which is the regime in which tracking is a question rather than a race against
the sampling period.

\subsubsection{Reference trajectory}
\label{subsubsec:raceline}
All four runs follow the same optimised raceline (Fig.~\ref{fig:raceline}),
generated with the TUM global race trajectory optimisation
package~\cite{tumracetraj} under its minimum-curvature
formulation~\cite{heilmeier2020}: a closed \SI{5828.2}{\meter} lap of
\num{2916} waypoints, reference speed
\SIrange{9.62}{25.78}{\meter\per\second}, minimum corner radius
\SI{18.7}{\meter} ($|\kappa|_{\max}=\SI{0.0536}{\per\meter}$). Its peak lateral
demand is $\max(v_{x}^{2}|\kappa|)=\SI{4.98}{\meter\per\second\squared}$
(\SI{0.51}{\g}), roughly \SI{73}{\percent} of the \SI{0.70}{\g} the static
surface sustains. On the decaying surface that same trajectory stops being
admissible part way through the run, which is the point of running it: the
plant passes below the reference's demand at an effective $\mu$ of about
\num{0.51}, and what the two arms do about it is the result of
Section~\ref{subsec:resultsclosedloop}.

\subsubsection{The two configurations}
\label{subsubsec:scenarios}
The two configurations of the pipeline differ in the nine settings of
Table~\ref{tab:scenarios} and in nothing else; vehicle, map, trajectory,
controller, data budget, loop count, tire prior and box constraints are shared,
and each configuration meets both surfaces with the same parameter file.
\basearm{} is the inherited configuration evaluated at full scale on the
corrected rollout speed, so the comparison isolates those nine settings rather
than reproducing the published pipeline in full; \oursarm{} is the shipped
configuration. Since the settings move together, what follows is a
configuration comparison, not a per-factor ablation, which
Section~\ref{subsec:ablations} supplies offline for the four that carry the
result. Nor is it a test of this paper's central contribution: the rollout
speed is taken from the measured lap in \emph{both} arms, and no run's mean
speed falls below \SI{16}{\meter\per\second}, so \oursarm{}'s
\SI{15}{\meter\per\second} floor never binds. The excitation correction itself,
together with the operating point it ships, is established in
Section~\ref{subsec:sweep}.

\begin{table}[t]
\caption{The two configurations. Each is run on both surfaces.}
\label{tab:scenarios}
\centering
\begin{threeparttable}
\footnotesize
\setlength{\tabcolsep}{4pt}
\begin{tabular}{@{}lll@{}}
\toprule
Setting & \basearm{} & \oursarm{}\\
\midrule
Residual architecture & MLP (\texttt{baseline}) & S4D (\texttt{s4})\\
Residual objective & One-step MSE & Physics-informed,
\eqref{eq:pinnloss}\\
Model integrator & Euler & RK4\\
Sweep amplitude & Fixed \SI{0.40}{\radian} & $1.5P_{99}(|\delta|)$, clipped\\
Pacejka solver & Trust-region~\cite{branch1999trust} & Differential
evolution~\cite{storn1997de}\\
Tikhonov weight $\lambda$ & \num{0} & \num{0.05}, frozen target\\
Friction warm start & Off & On; $D$ pinned per axle\\
Fast-$\mu$ derate & Off & On, \SI{1}{\hertz} on \SI{6}{\second}\\
Update relaxation $\beta$ & \num{1.00} (undamped) & \num{0.20}\\
\bottomrule
\end{tabular}
\begin{tablenotes}[flushleft]\scriptsize
\item Everything not listed is identical between the two, the tire prior and
the box constraints included. $\beta$ is the relaxation of \eqref{eq:relax};
\num{1.00} is the inherited undamped update. The fast-$\mu$ derate is the brush
estimator of Section~\ref{subsec:friction} refitted on a trailing
\SI{6}{\second} window between coefficient cycles and published as an optional
ratio the controller applies to the identified grip ceiling, clamped to
$[\num{0.4},\num{1.2}]$; with the warm start off there is nothing to publish.
The rollout speed is taken from the measured lap in both arms and is not a
difference here.
\end{tablenotes}
\end{threeparttable}
\end{table}

\subsubsection{Provenance of the tire prior}
\label{subsubsec:priorprovenance}
The prior $\phip^{0}$ both arms start from, and that the Tikhonov term of
\eqref{eq:fitobj} pulls toward, is the measured tire of Table~\ref{tab:ref},
obtained by fitting \eqref{eq:mf} to CARLA's own per-wheel telemetry over a
dedicated limit-excitation run that reached $\alpha=\SI{0.37}{\radian}$, axle
$\mu=\num{1.05}$ and \SI{1.50}{\g}. It is not independent of the validation
data, and the manoeuvre that produced it is one the pipeline is not allowed to
run, the raceline peaking at \SI{0.51}{\g}. Decisively, it was measured on a
surface none of these runs use, carrying $D=\num{1.009}/\num{1.002}$ and axle
slopes of \num{13.2} and \SI{13.9}{\kilo\newton\per\radian} against the
\num{0.70} and the \num{17.4} and \SI{16.6}{\kilo\newton\per\radian} the static
surface realises. Returning it unchanged is a \SI{44}{\percent} grip
overstatement on the static surface and a \SI{188}{\percent} one at the
decaying surface's floor, so an arm has to move away from that prior, in the
right direction, to score well, and on the decaying surface it has to keep
moving.

\subsection{Selecting the rollout operating point $(v,\delta_{\mathrm{sweep}})$}
\label{subsec:sweep}

This grid is where the rollout configuration the pipeline ships is
\emph{chosen}, not an auxiliary check on a choice made elsewhere: the
\SI{15}{\meter\per\second} speed floor and the
$\delta_{\max}=\SI{0.54}{\radian}$ clip of Section~\ref{subsec:excitation} are
read off Table~\ref{tab:sweep}, and the reachable-friction bound
\eqref{eq:mureach} is what makes the search one-dimensional rather than a blind
scan. The search is run where the answer is known exactly rather than in CARLA,
since a CARLA lap cannot be scored against a $D$ nobody knows. A tire with a
known peak factor $D$ is rolled out through the synthetic sweep of stage~2 (the
shipped rollout, residual network set to zero) and fitted back with the shipped
solver, \emph{unregularised} and started from the true coefficients, so
everything it loses is lost to excitation. Rollout speed $v$, sweep amplitude
$\delta_{\mathrm{sweep}}$ and true $D$ are swept, each cell repeated over five
noise realisations at the one-step persistence RMSE measured on CARLA runs of
this vehicle. The configuration axes are swept one at a time, because
$\mu_{\mathrm{reach}}$ is a function of both. The selection criterion is the one
the pipeline needs: the widest range of true $D$ recovered without railing, over
the whole \numrange{0.70}{1.80} column span, since the tire is unknown when the
configuration is fixed. The span matters twice over on a decaying surface, where
one configuration has to cover every $D$ the run passes through.

\begin{table*}[t]
\caption{The search that fixes the shipped rollout operating point: peak factor
$\hat D$ recovered from a known tire, swept along both arguments of
$\mu_{\mathrm{reach}}$. The shipped $(v,\delta_{\mathrm{sweep}})$ is selected
from this grid.}
\label{tab:sweep}
\centering
\begin{threeparttable}
\footnotesize
\setlength{\tabcolsep}{6pt}
\begin{tabular}{@{}rrcccc@{}}
\toprule
swept & $\mu_{\mathrm{reach}}$ & \multicolumn{4}{c}{$\hat D$ (railed of 8) at true $D$ of}\\
\cmidrule(l){3-6}
value & & 0.70 & 1.05 & 1.40 & 1.80\\
\midrule
\multicolumn{6}{@{}l}{\emph{(a) rollout speed $v$ [\si{\meter\per\second}] swept at $\delta_{\mathrm{sweep}}=\SI{0.40}{\radian}$}}\\
4 & 0.43 & \num{0.52}$\pm$\num{0.036} (1.8) & \num{0.57}$\pm$\num{0.013} (2.8) & \num{0.59}$\pm$\num{0.018} (3.8) & \num{0.50}$\pm$\num{0.001} (7)\\
6 & 0.96 & \num{0.68}$\pm$\num{0.001} (1) & \num{1.04}$\pm$\num{0.014} (1.2) & \num{1.22}$\pm$\num{0.066} (1.8) & \num{1.26}$\pm$\num{0.097} (2)\\
8 & 1.70 & \num{0.69}$\pm$\num{0.001} (1) & \num{1.02}$\pm$\num{0.001} (0.8) & \num{1.36}$\pm$\num{0.002} (0) & \num{1.79}$\pm$\num{0.025} (1.2)\\
13 & 4.49 & \num{0.72}$\pm$\num{0.001} (0) & \num{1.05}$\pm$\num{0.002} (0) & \num{1.38}$\pm$\num{0.001} (0.4) & \num{1.77}$\pm$\num{0.002} (0.4)\\
20 & 10.64 & \num{0.92}$\pm$\num{0.004} (1) & \num{1.14}$\pm$\num{0.005} (0.4) & \num{1.43}$\pm$\num{0.002} (0) & \num{1.80}$\pm$\num{0.005} (0)\\
\addlinespace
\multicolumn{6}{@{}l}{\emph{(b) $\delta_{\mathrm{sweep}}$ [\si{\radian}] swept at $v=\SI{15}{\meter\per\second}$}}\\
0.03 & 0.45 & \num{0.80}$\pm$\num{0.019} (2) & \num{1.10}$\pm$\num{0.041} (2) & \num{1.08}$\pm$\num{0.182} (2) & \num{1.09}$\pm$\num{0.397} (2)\\
0.06 & 0.90 & \num{0.92}$\pm$\num{0.003} (2) & \num{1.30}$\pm$\num{0.008} (2) & \num{1.68}$\pm$\num{0.017} (2) & \num{1.80}$\pm$\num{0.232} (3.4)\\
0.15 & 2.24 & \num{0.99}$\pm$\num{0.004} (2) & \num{1.38}$\pm$\num{0.003} (2) & \num{1.71}$\pm$\num{0.019} (2) & \num{1.98}$\pm$\num{0.009} (1)\\
0.40 & 5.98 & \num{0.76}$\pm$\num{0.009} (1) & \num{1.05}$\pm$\num{0.002} (0) & \num{1.39}$\pm$\num{0.002} (0) & \num{1.78}$\pm$\num{0.001} (0)\\
0.54 & 8.08 & \num{0.83}$\pm$\num{0.004} (1) & \num{1.08}$\pm$\num{0.002} (0) & \num{1.38}$\pm$\num{0.001} (0) & \num{1.78}$\pm$\num{0.008} (0)\\
\bottomrule
\end{tabular}
\begin{tablenotes}[flushleft]\scriptsize
\item Unregularised, started from the true coefficients. Each cell is $\hat D$ $\pm$ the half-width of its \SI{95}{\percent} $t$-interval over five noise realisations, with the mean number of the eight coefficients resting on a box constraint in parentheses. The two axles of one rollout are not independent, so each realisation contributes the mean of its own $(\hat D_{f},\hat D_{r})$ and $n=5$.
\item $\mu_{\mathrm{reach}}$ is \eqref{eq:mureach} at $L=\SI{1.5334}{\meter}$; the inherited configuration is the first row of (a). It is the Ackermann \emph{upper} bound and above $\mu_{\mathrm{reach}}\approx2$ detaches from what the rollout realises, so a large value means ``far past saturation'', not a demand the tire ever sees.
\end{tablenotes}
\end{threeparttable}
\end{table*}

The grid says three things. First, the inherited \SI{4}{\meter\per\second}
rollout is uninformative about $D$ at \emph{every} grip level tested: the
$0.70$, $1.05$, $1.40$ and $1.80$ tires all come back as $\hat D$ between
\num{0.50} and \num{0.59}, with two to seven of the eight coefficients on a
bound, so the fit reports $\mu_{\mathrm{reach}}$ and not the tire. Second, the
boundary is $\mu_{\mathrm{reach}}\gtrsim D$ and not a speed: at
\SI{6}{\meter\per\second}, where $\mu_{\mathrm{reach}}=0.96$, the $0.70$ and
$1.05$ tires are recovered to within \SI{3}{\percent} while the $1.40$ and
$1.80$ tires are \SI{13}{\percent} and \SI{30}{\percent} short, and for
$\mu_{\mathrm{reach}}$ between \num{1.7} and \num{4.5} every tire is recovered
to within \SI{4}{\percent} with no systematic railing. Nor is the bound a
quantity to maximise: at $v=\SI{20}{\meter\per\second}$,
$\mu_{\mathrm{reach}}=\num{10.64}$, the rollout sits far past saturation for
every tire on the grid and $\hat D$ re-biases upward, the \num{0.70} tire
returning \num{0.92}. Third, the failure is not the solver's: multi-start, a
global solver or a better start point cannot recover a coefficient the data
does not contain.

Panel (b) sweeps the other argument of \eqref{eq:mureach}. Starving the
amplitude destroys $D$ just as starving the speed does, so speed is not the
governing variable, and the rollout rather than the tire sets the direction of
the error: an under-speeded rollout returns \emph{too little} grip, an
under-swept one too much (at $\delta_{\mathrm{sweep}}=\SI{0.06}{\radian}$ the
\num{1.05} tire comes back at \num{1.30} and the \num{1.40} tire at \num{1.68},
two of the eight coefficients on a bound throughout those rows). That is the
signature of an unidentified parameter, not a biased one. The axes do
\emph{not} collapse onto a common $\mu_{\mathrm{reach}}$, for the upper-bound
reason of Section~\ref{subsec:excitation}: at $v=\SI{20}{\meter\per\second}$
\eqref{eq:mureach} reads \num{10.64} against a realised
\numrange{1.40}{2.18}. Read against the realised demand both axes behave alike,
so \eqref{eq:mureach} is a cheap screening bound, conclusive when it falls
below $D$, and not a measure of excitation achieved.

The two panels fix the shipped point between a lower edge and an upper one. On
speed, railing loses the \num{1.40} and \num{1.80} tires at
$v\le\SI{6}{\meter\per\second}$ and the saturation re-bias spoils the
\num{0.70} tire at \SI{20}{\meter\per\second}; only $v=\num{8}$ and
$v=\num{13}$ recover every tire in the span to within \SI{4}{\percent}, and the
\SI{15}{\meter\per\second} floor sits inside that band with the measured lap
speed governing above it. On amplitude,
$\delta_{\mathrm{sweep}}\le\SI{0.15}{\radian}$ holds two coefficients on a
bound at every $D$, while \SI{0.40}{\radian} and \SI{0.54}{\radian} rail
nothing, so \SI{0.54}{\radian} is taken as $\delta_{\max}$, the ceiling of the
$1.5P_{99}$ clip that selects the amplitude on a given lap. All of it is a best
case: with the residual set to zero a Magic Formula is fitted to Magic-Formula
data under exactly matched structure, so the selected point is the best setting
\emph{on this grid} and not a proven optimum.

\subsection{Cross-run summary}
\label{subsec:summary}

Table~\ref{tab:comparison} is the complete comparison over both surfaces. On
the static surface \oursarm{} is better on every quantity scored against the
plant and behind only on QP solve time, by \SI{5}{\percent} of the RMS and
\SI{1}{\percent} of the worst case, both two orders of magnitude inside the
control period. The margins are \numrange{9.0}{12.1}$\times$ on peak friction
and axle force and \num{1.5}$\times$ on MPC tracking error.

On the decaying surface the comparison stops being a margin and becomes an
outcome. \oursarm{} completes the same three timed laps, at
\SI{0.139}{\meter} RMS lateral error under MPC, which is its static-surface
figure to within \SI{8}{\percent} although half the grip is gone by the end.
\basearm{} leaves the track in lap~2 and the run is scored as a failure, its
\SI{23.7}{\meter} RMS error under MPC being the departure itself rather than a
tracking quality. The peak-friction column is where the two diverge first:
\oursarm{} follows the surface down, \basearm{} publishes $D=\num{2.000}$
throughout, and the \num{1.408} error that implies is what the controller is
holding when the trajectory stops being admissible. The rest of this section
reads those rows against the figures; Section~\ref{subsec:ablations} attributes
them.

\begin{table*}[t]
\caption{Cross-run comparison. Two configurations on two surfaces.}
\label{tab:comparison}
\centering
\begin{threeparttable}
\footnotesize
\setlength{\tabcolsep}{5pt}
\begin{tabular}{@{}lrrrr@{}}
\toprule
& \multicolumn{2}{c}{Static surface, $\mu=0.700$} & \multicolumn{2}{c}{Decaying surface, $\mu:0.700\rightarrow0.350$}\\
\cmidrule(lr){2-3}\cmidrule(l){4-5}
Metric & \oursarm{} & \basearm{} & \oursarm{} & \basearm{}\\
\midrule
\multicolumn{5}{@{}l}{\emph{Outcome and closed-loop tracking (wall clock)}}\\
Timed laps completed & \textbf{3} & \textbf{3} & \textbf{3} & 1 (off track, lap 2)\\
Run duration [\si{\second}] & 580.4 & 554.6 & 640.9 & 358.7\\
Lap times [\si{\second}] & 175.2 / 145.4 / 142.8 & 162.3 / 135.8 / 138.0 & 189.8 / 160.9 / 169.2 & 167.7 / --- / ---\\
RMS $e_{y}$, overall [\si{\meter}] & \textbf{0.256} & 0.319 & \textbf{0.264} & 21.505\\
RMS $e_{y}$, MPC [\si{\meter}] & \textbf{0.151} & 0.224 & \textbf{0.139} & 23.705\\
Max $|e_{y}|$, MPC [\si{\meter}] & \textbf{0.683} & 0.776 & \textbf{0.664} & 38.386\\
ITAE $e_{y}$, MPC [\si{\meter\second\squared}] & \textbf{15925} & 25384 & \textbf{17888} & $1.32\times10^{6}$\\
MPC solve, RMS / max [\si{\milli\second}] & 0.399 / 2.604 & \textbf{0.380 / 2.568} & \textbf{0.420 / 3.589} & 0.678 / 4.062\\
\addlinespace
\multicolumn{5}{@{}l}{\emph{Axle lateral force against per-wheel telemetry (simulation clock)}}\\
Front / rear RMSE [\si{\newton}] & \textbf{27.2 / 33.2} & 329.6 / 348.4 & \textbf{32.9 / 36.5} & 171.4 / 178.3\\
Front / rear $R^{2}$ & \textbf{0.993 / 0.987} & $-0.160$ / $-0.461$ & \textbf{0.984 / 0.977} & 0.708 / 0.647\\
Front / rear max abs.\ error [\si{\newton}] & \textbf{264 / 135} & 1295 / 1316 & \textbf{249 / 236} & 2507 / 2069\\
\addlinespace
\multicolumn{5}{@{}l}{\emph{Peak friction against the plant's effective $\mu$}}\\
Front / rear $\mu$ RMSE, whole run & \textbf{0.145 / 0.142} & 1.300 / 1.300 & \textbf{0.209 / 0.207} & 1.408 / 1.408\\
Front / rear $\mu$ RMSE, after release & \textbf{0.078 / 0.078} & 1.300 / 1.300 & \textbf{0.149 / 0.149} & 1.408 / 1.408\\
\addlinespace
\multicolumn{5}{@{}l}{\emph{One-step state prediction}}\\
$v_{y}$ / $\omega$ RMSE [\si{\meter\per\second}, \si{\radian\per\second}] & \textbf{0.014 / 0.049} & 0.156 / 0.104 & \textbf{0.011 / 0.037} & 0.055 / 0.105\\
\addlinespace
\multicolumn{5}{@{}l}{\emph{Last accepted coefficient set, front / rear}}\\
$B$ & 12.89 / 13.75 & 7.64 / 11.19 & 16.10 / 16.42 & \textbf{4.00} / \textbf{4.00}\\
$C$ & 1.240 / 1.248 & 1.806 / 1.346 & 1.362 / 1.416 & 1.204 / 1.591\\
$D$ & 0.768 / 0.768 & \textbf{2.000} / \textbf{2.000} & 0.460 / 0.460 & \textbf{2.000} / \textbf{2.000}\\
$E$ & $-2.52$ / $-1.84$ & \textbf{$-3.00$} / $-1.77$ & $-2.08$ / $-1.40$ & \textbf{$-3.00$} / $-1.74$\\
Axle stiffness [\si{\kilo\newton\per\radian}] & 16.9 / 16.8 & 37.9 / 38.4 & 13.8 / 13.6 & 13.2 / 16.2\\
Grip ceiling \eqref{eq:gripceiling} [\si{\g}] & 0.77 & 2.00 & 0.46 & 2.00\\
Plant grip at that moment [\si{\g}] & 0.70 & 0.70 & 0.35 & 0.53\\
\bottomrule
\end{tabular}
\begin{tablenotes}[flushleft]\scriptsize
\item Bold marks the better arm within a surface for the metric blocks, and a
coefficient resting on a box constraint in the coefficient block. The box is
$\phip\in([4.0,1.2,0.4,-3.0],[20.0,2.2,2.0,1.0])$ per axle. The \basearm{}
decaying-surface row is the last set it published before leaving the track, and
its tracking figures are the departure, not a tracking quality. Lap times run on
the ROS wall clock while the simulator advances its fixed step faster than real
time, at a ratio that differs between runs, so they are not read as a
vehicle-performance comparison across arms; within the \oursarm{} column the
static-to-decaying difference is discussed in
Section~\ref{subsec:resultsclosedloop}.
\end{tablenotes}
\end{threeparttable}
\end{table*}

\subsection{Identified tire model}
\label{subsec:resultstire}

The coefficient block of Table~\ref{tab:comparison} separates the two arms
before any downstream metric does.

First, no coefficient of \oursarm{} rests on a box constraint in any accepted
identification of either run. \basearm{} puts $D$ on its upper bound of
\num{2.000} on both axles in \emph{every} accepted identification of both runs
and $E_{f}$ on its lower bound of \num{-3.0} likewise, and on the decaying
surface it adds $B$ on its lower bound of \num{4.0} on both axles, so five of
its eight coefficients are resting on an edge in the set the car was driving on
when it left the track. The \emph{inner} fits are a separate matter:
\oursarm{}'s static run logs \num{35} fits that stopped on a bound, \num{30} of
them on $C$'s lower \num{1.2} and none on $D$, which the pin takes out of the
fit, against \num{96} for \basearm{} spread over $B$, $C$ and $D$; on the
decaying surface the count rises to \num{56}, still concentrated on $C$. It is the
relaxation of \eqref{eq:relax} that keeps the accepted set off an edge its
iterates touch. A peak factor of \num{2.000} against a plant realising
\num{0.700} is less a \SI{186}{\percent} parameter error than the box-edge
signature of Section~\ref{subsec:identifiability}: the coefficient is not
identified, and no per-coefficient bounds check distinguishes that from a fit.

Second, the linear range is identified rather than inherited. Every stiffness in
the table is \eqref{eq:stiffness} at the static axle loads
$F_{z,f}=\SI{1373}{\newton}$, $F_{z,r}=\SI{1274}{\newton}$; for the plant it is
the slope fitted to the simulator's own per-wheel forces over
$|\alpha|\in[\num{0.005},\num{0.02}]\,\si{\radian}$, so the two are comparable.
The static surface realises \SI{17.4}{\kilo\newton\per\radian} front and
\SI{16.6}{\kilo\newton\per\radian} rear over this run's load range, below the
$\bar F_{z}=1$ closed-form values of Table~\ref{tab:ref} for the reason its
first footnote gives. The prior carries \num{13.2} and \num{13.9},
\SI{24}{\percent} and \SI{16}{\percent} low; \oursarm{} returns \num{16.9} and
\num{16.8}, \SI{3.2}{\percent} low front and \SI{1.4}{\percent} high rear,
removing \SI{87}{\percent} and \SI{91}{\percent} of the prior's error, and the
product $BCDF_{z}$ that the rollout fit determines
(Section~\ref{subsec:bcd}) then holds within
\SIrange{16.0}{16.9}{\kilo\newton\per\radian} over the five cycles after the
gate releases. \basearm{} returns \num{37.9} and \num{38.4},
\SI{118}{\percent} and \SI{132}{\percent} high, which is $D=\num{2.0}$
multiplying a $B$ the fit chose against it. On the decaying surface
\oursarm{}'s stiffness falls with the surface, to \num{13.8} and
\SI{13.6}{\kilo\newton\per\radian} at the last accepted set, while \basearm{}
arrives at \num{13.2} and \SI{16.2}{\kilo\newton\per\radian} by a different
route: $B$ at its floor multiplied by $D$ at its ceiling. A plausible stiffness
assembled from two railed coefficients is the clearest case this paper has for
gating on derived quantities rather than on coefficients, since the stiffness
gate passes it.

Third, both sets keep the understeer sign and pass the stability gate of
Section~\ref{subsec:gates}; what separates them is the grip ceiling, which
\eqref{eq:gripceiling} puts at \SI{2.00}{\g} for \basearm{} on both surfaces
against \SI{0.77}{\g} and \SI{0.46}{\g} for \oursarm{}. The baseline's ceiling
is a claim the plant cannot honour by a factor of \num{2.8} on the static
surface and \num{3.8} at the moment the car leaves the decaying one, and the
grip gate is one-sided, so it passes both. All of these readings are possible
because the lap traverses part of the curvature:
$P_{99}(|\alpha|)=\SI{2.7}{\degree}$ and $\max|\alpha|=\SI{5.8}{\degree}$
against the \SI{8.0}{\degree} at which the plant's curve peaks on the static
surface, and $\max|F_{y}|/F_{z}$ reaches \num{0.71} front and \num{0.72} rear
in transients against \num{0.49} and \num{0.50} at $P_{99}$. This is not a
limit manoeuvre, but it suffices for the friction gate to release, which on the
$\mu=\num{1.05}$ runs it did not; as the decaying surface falls the same lap
becomes relatively more demanding, so the gate keeps releasing rather than
closing.

\subsection{Peak friction: holding one surface, tracking another}
\label{subsec:resultsmu}

Figure~\ref{fig:mutimeseries} shows what each arm published for the peak factor
over both runs, and Table~\ref{tab:comparison} gives the errors behind it. The
pinned-$D$ design of Section~\ref{subsec:friction} is directly observable here,
because both arms start \SI{44}{\percent} away from the static answer and
\SI{188}{\percent} away from the decaying one's floor.

On the static surface \oursarm{} publishes six accepted identifications. The
first carries $D_{f}=\num{1.009}$, $D_{r}=\num{1.002}$: the friction gate
rejected that buffer, so the axles held the prior they came in with, the
hold-rather-than-fit rule of Section~\ref{subsec:frictionfloor} and not a fitted
value. From $t=\SI{140.6}{\second}$ the gate releases on every subsequent cycle
and the published peak factor drops to \num{0.779}, \num{0.821}, \num{0.749},
\num{0.756} and \num{0.768}, front and rear alike, against a plant at
\num{0.700}: \SI{10.7}{\percent} high on average, from a prior that was
\SI{44}{\percent} high. The whole-run $\mu$ RMSE is dominated by the first
\SI{140}{\second} at the prior; after the first release it is \num{0.078} on
both axles.

On the decaying surface the same six cycles land on a surface that has moved
between each of them. The gate again rejects the first buffer and releases from
$t=\SI{145.1}{\second}$, and the published peak factor then falls monotonically:
\num{0.706}, \num{0.679}, \num{0.609}, \num{0.525} and \num{0.460}, against a
plant falling from \num{0.580} to \num{0.350} over the same interval. Read as a
rate, the identifier recovers the decay itself: the published sequence falls at
$4.28\times10^{-4}\,\si{\per\second}$ against the plant's
$4.00\times10^{-4}\,\si{\per\second}$, a
\SI{7}{\percent} error on the quantity that a static surface cannot test at all.
Read as a level it carries two offsets, and separating them matters. Each cycle
is fitted on the buffer that preceded it, so the published value describes the
surface of roughly one cycle earlier; against the plant one cycle back the
residual is \SIrange{10.9}{17.4}{\percent} high, mean \SI{14.3}{\percent},
against the \SI{10.7}{\percent} the same estimator carries on a surface that
does not move. The apparent whole-run RMSE of \num{0.209} is therefore that
familiar brush-model bias plus one cycle of transport lag on a surface losing
about \num{0.06} of $\mu$ per cycle, and not a new error mode. The remedy for
the lag is a shorter cycle rather than a better fit, and the fast-$\mu$ derate
of Table~\ref{tab:scenarios} is that shorter cycle; on this run it spent the
last cycles at the upper end of its clamp and so did not shorten anything
(Section~\ref{subsec:ablations}).

\basearm{} returns $D=\num{2.000}$ on both axles in every accepted
identification of both runs and never leaves the bound. On the static surface
that is a constant $\mu$ error of \num{1.300}, so \oursarm{} is
\num{9.0}$\times$ better on the whole-run figure and \num{16.7}$\times$ better
once its warm start has released. On the decaying surface the error grows as
the plant walks away from the bound, reaching \num{1.408} RMSE before the run
ends early, and \oursarm{} is \num{6.7}$\times$ better there. Two
qualifications belong on this whole subsection: the peak factor is
\emph{measured} by the brush estimator on the recorded trace and then held, not
recovered by the Pacejka fit, which Section~\ref{subsec:bcd} shows cannot
recover it; and the residual bias is one-signed, consistent with the brush
model's tendency to over-read $\mu$ where the curve is not fully traversed,
which is also why the tracking result below has margin to work with.
Appendix~\ref{subsec:mu} evaluates the same estimator on a synthetic brush-tire
plant, where $\mu$ is known exactly and the excitation can be swept
independently of it, and shows that it is flat in excitation above its gate
while the utilisation bound it replaces is not.

\begin{figure*}[t]
\centering
% Front-axle column only: the rear panels are numerically identical (D is pinned
% to the same value on both axles), so they are cropped to give the front panels
% the full width. Trim is in PDF points on the 864x460.8 pt page shared by both
% files; the gutter between the two panel columns sits at x = 652..669 px of a
% 1320 px raster, and the top 40 px carry the now-redundant suptitle. The bottom
% 25 pt are the strip Comparison.save reserves in both files for the
% incomplete-run banner, which the right trim would otherwise slice mid-word;
% the caption states the same fact.
\includegraphics[trim=0 25 431 26, clip, width=0.49\textwidth]{\cmpstatic/mu_timeseries_by_scenario}\hfill
\includegraphics[trim=0 25 431 26, clip, width=0.49\textwidth]{\cmpdecay/mu_timeseries_by_scenario}
\caption{Peak friction over the run, front axle, one row per configuration:
the plant's effective $\mu$ (black) against the identified peak factor $D$
published at each accepted identification. Left, the static surface; right, the
decaying surface. The rear axle is identical and is omitted, since $D$ is pinned
to the same value on both axles (Section~\ref{subsec:frictionfloor}). On the
static surface every departure from the black line is estimator error:
\oursarm{} publishes the $\mu=\num{1.05}$ prior while the gate is closed, steps
down onto the plant within \SI{11}{\percent} at the first release, and holds.
On the decaying surface the same five accepted cycles walk the published $D$
down at \SI{7}{\percent} of the plant's own rate, offset by one cycle of
transport lag and the same estimator bias. \basearm{} sits on the
$D\le\num{2.0}$ box constraint throughout both runs, and its right-hand panel
ends where the car left the track.}
\label{fig:mutimeseries}
\end{figure*}

\subsection{Axle lateral force and one-step state prediction}
\label{subsec:resultsforce}

The identified curve is scored at the measured $(\alpha,F_{z})$ against the
plant's own lateral force, summed per axle (Table~\ref{tab:comparison}). On the
static surface \oursarm{} leads on both axles by an order of magnitude in RMSE
and in the error tails; the baseline's negative $R^{2}$ means its curve predicts
the plant's axle force worse than the mean of that force does. This is the
linear-range result seen as force: at $P_{99}(|\alpha|)=\SI{2.7}{\degree}$ the
metric is dominated by cornering stiffness, so the models separate almost
entirely on slope. Because $D$ multiplies that slope in \eqref{eq:stiffness},
\basearm{}'s railed $D$ is the direct cause of its force error. Pinning $D$
alone does not produce these numbers: with $\beta=1$ the same arm degrades
monotonically with iteration count (Section~\ref{subsec:ablations}), which is
why the relaxed update of \eqref{eq:relax} belongs in the method rather than a
configuration file.

On the decaying surface \oursarm{}'s force error grows by about a fifth,
\SI{32.9}{\newton} front against \SI{27.2}{\newton}, which is the transport lag
of Section~\ref{subsec:resultsmu} converted into force, and $R^{2}$ stays at
\num{0.984} and \num{0.977}. \basearm{}'s force error is \emph{lower} there than
on the static surface, \SI{171.4}{\newton} against \SI{329.6}{\newton}, and its
$R^{2}$ turns positive. That is not a success: its run is scored over \num{2597}
samples against \num{10744}, it ends before the surface reaches its floor, and
its worst-case error is the largest in the table at \SI{2507}{\newton}. A model
with five coefficients on a box edge produced a better-looking aggregate over a
shorter and slower run, which is the diagnostic problem
Section~\ref{subsec:identifiability} describes.

Level~1 of the validation protocol is reported against a persistence baseline
(next state equals last measured state) written by the same benchmark node.
\oursarm{} leads on both signals on both surfaces, by an order of magnitude on
lateral velocity on the static surface ($R^{2}$ \num{0.998} against
\num{0.709}) and by a factor of two on yaw rate ($R^{2}$ \num{0.892} against
\num{0.527}). Neither model beats persistence at the \SI{0.035}{\second} model
step for either configuration or signal, because a one-step metric on a smooth
trajectory rewards not moving. One-step accuracy is therefore weak evidence for
either tire model, and the gap is read only as confirming the direction the
axle-force and closed-loop comparisons establish.

\subsection{Closed-loop tracking}
\label{subsec:resultsclosedloop}

Figure~\ref{fig:overlay} puts the two surfaces side by side over the whole of
each run.

On the static surface both arms finish three timed laps with zero emergency
halts and the same finite-state sequence, with the handover at
$t=\SI{79.1}{\second}$ for \oursarm{} and $t=\SI{64.0}{\second}$ for
\basearm{}. Under MPC, where the identified model enters the loop, \oursarm{}
tracks at \SI{0.151}{\meter} RMS lateral error against \SI{0.224}{\meter}, a
\SI{32}{\percent} reduction, with maximum \SI{0.683}{\meter} against
\SI{0.776}{\meter}. ITAE is $\int t|e_{y}|\,\mathrm{d}t$, neither scale-free
nor comparable across runs of different length, so the duration-normalised
forms $\mathrm{IAE}/T=\num{0.104}$ against \num{0.157} and
$\mathrm{ITAE}/(T^{2}/2)=\num{0.095}$ against \num{0.166} are reported beside
it; all three order the runs alike. The split by controller confirms that the
whole difference sits in the MPC phase: the pure-pursuit phases, which do not
use the identified model, give \num{110} and \num{107}, within
\SI{3}{\percent} of each other. The QP cost is unchanged, both arms staying two
orders of magnitude below the \SI{50}{\milli\second} control period; the
identified coefficients change the model the QP carries, not its size.

On the decaying surface the two arms separate at $t\approx\SI{246}{\second}$,
in lap~2. \basearm{}'s lateral error steps from under \SI{1}{\meter} to
\SI{38.4}{\meter} inside a few seconds, its speed goes through zero to
$-16.9\,\si{\meter\per\second}$, and the run ends on a collision whose horizontal
impulse of \SI{415.9}{\newton\second} is twice the \SI{200}{\newton\second}
abort threshold. The plant's effective $\mu$ at that moment is \num{0.525},
worth \SI{0.53}{\g}, against a reference that is still asking for
\SI{0.51}{\g} at its peak and a model that is still reporting a
\SI{2.00}{\g} ceiling. The failure is not a tracking degradation that grew
until it mattered; the tracking error is flat and small right up to the
departure, which is what a grip limit being crossed looks like from inside the
controller.

\oursarm{} completes its three laps on the same surface and its tracking error
does not degrade as the grip halves: binned in \SI{100}{\second} windows after
the handover it reads \num{0.153}, \num{0.138}, \num{0.155}, \num{0.094},
\num{0.152} and \SI{0.092}{\meter} RMS, with no trend, against \num{0.123} to
\SI{0.179}{\meter} over the corresponding windows of the static-surface run.
What moves instead is the speed. The identified grip ceiling falls with the published $D$, the
controller re-clamps the reference speed profile against it, and the peak
lateral demand of the clamped reference drops from \SI{4.98}{\meter\per\second\squared}
(\SI{0.51}{\g}) at the start of the run to \SI{2.92}{\meter\per\second\squared}
(\SI{0.30}{\g}) at the end. The three lap times lengthen accordingly, to
\num{189.8}, \num{160.9} and \SI{169.2}{\second} against \num{175.2},
\num{145.4} and \SI{142.8}{\second} on the static surface, an \SI{18}{\percent}
loss on the last lap. That trade is the whole argument for identifying on
track: the vehicle pays lap time it has to pay, and the identification is what
tells it how much. \basearm{} never makes the trade, because a peak factor
pinned to a box constraint cannot tell the controller that anything has
changed. Lap times themselves are not read as a vehicle-performance result
across arms, for the clock reason of Section~\ref{subsec:setup}; the comparison
that carries this paragraph is within the \oursarm{} column, where the same
configuration meets two surfaces.

\begin{figure*}[t]
\centering
\includegraphics[width=0.49\textwidth]{\cmpstatic/control_timeseries_overlay_by_scenario}\hfill
\includegraphics[width=0.49\textwidth]{\cmpdecay/control_timeseries_overlay_by_scenario}
\caption{Lateral error, longitudinal speed and MPC solve time over the whole of
each run, both configurations overlaid, shaded by the active controller state
of the \oursarm{} arm. Left, the static surface, where both arms run to three
laps and the traces stay superimposed. Right, the decaying surface, where
\basearm{} departs the track at $t\approx\SI{246}{\second}$ ($e_{y}$ to
\SI{38.4}{\meter}, speed through zero) while \oursarm{} runs to
\SI{640.9}{\second}. \oursarm{}'s speed trace on the right loses its peaks over
the second half of the run: the controller re-clamps the reference against a
grip ceiling that is falling because the identifier says it is. Solve time
stays two orders of magnitude below the \SI{50}{\milli\second} control period
in every case.}
\label{fig:overlay}
\end{figure*}

\subsection{What carries the result}
\label{subsec:ablations}

The two arms differ in nine settings at once, so the comparison above
establishes an ordering and attributes nothing. Appendix~\ref{sec:ablations}
attributes it by offline replay of the identification stage on the five
\SI{30}{\second} buffers each arm collected during its own run, scored with the
quantity the closed-loop tire-force benchmark reports. Four settings carry the
result. Refusing to fit an unmeasured $D$ removes most of the error, cutting the
mean $|D-\mu|$ from \num{0.585} to \num{0.084} on replayed buffers of a
$\mu=\num{1.05}$ run. Damping the loop carries the rest: axle force RMSE is
\num{70.4} and \SI{62.3}{\newton} at the inherited $\beta=\num{1.00}$,
\num{92.5} and \SI{71.8}{\newton} at $\beta=\num{0}$, which is no
identification at all, and \num{28.0} and \SI{19.4}{\newton} at the
$\beta=\num{0.20}$ that ships. Normalising the physics-informed objective is
necessary but not sufficient alone, and the S4D residual is worth
\SI{21}{\percent} on the force metric at matched settings. Two settings buy
nothing measurable: the residual \emph{objective}, indistinguishable from a
plain one-step MSE once the normalisation and the damping are both in place,
and the fast-$\mu$ derate, which sat at the upper end of its
$[\num{0.4},\num{1.2}]$ clamp over the last cycles of the decaying run, so the
reference re-clamp of Section~\ref{subsec:resultsclosedloop} is carried by the
published $D$ and not by the derate.

\subsection{Identification cost}
\label{subsec:runtime}

One S4D identification, six co-identification iterations on a \SI{30}{\second}
lap, takes \SI{20.9}{\second} end to end, inside the \SI{30}{\second}
re-identification period. The reduction from the \SI{181.4}{\second} it
originally took is exact: no epoch count, iteration count, early-stopping rule
or loss weight changed, and the identified coefficients agree to four decimal
places. That margin is what makes the decaying-surface result possible at all,
and it is also what bounds it: one cycle of lag on a surface moving at
\eqref{eq:decay}'s rate costs the \SIrange{10.9}{17.4}{\percent} offset of
Section~\ref{subsec:resultsmu}, and a surface moving ten times faster would cost
ten times that. The per-phase breakdown is in the supplementary material.
